% Captions for universal-partition-depth-observatory panels.
% Include these in the main .tex via % Captions for universal-partition-depth-observatory panels.
% Include these in the main .tex via % Captions for universal-partition-depth-observatory panels.
% Include these in the main .tex via % Captions for universal-partition-depth-observatory panels.
% Include these in the main .tex via \input{figures/universal-captions.tex}.

\begin{figure*}[!htbp]
    \centering
    \includegraphics[width=\textwidth]{figures/panel_1_composition.png}
    \caption{\textbf{Composition-inflation mechanism.}
    (\textbf{A})~Labelled composition count $T(n, 3) = 3\cdot 4^{n-1}$ versus oscillation depth $n$ on semilog axes. The straight line on the log plot confirms exponential growth with base $d+1 = 4$. Representative values: $T(10, 3) = 7.86\times 10^5$, $T(30, 3) = 3.89\times 10^{18}$, $T(50, 3) = 4.28\times 10^{30}$, demonstrating that even modest depths generate astronomical state counts.
    (\textbf{B})~Comparison of $T(n, d)$ across dimensions $d \in \{2, 3, 4, 5\}$ on semilog axes (viridis colouring). Higher $d$ produces steeper growth because the base $d+1$ increases; the slope ratio between $d=2$ ($3^{n-1}$) and $d=5$ ($6^{n-1}$) is $\log(6)/\log(3) \approx 1.63$. This establishes that dimension is a free parameter; the physical case $d=3$ is chosen for the three-coordinate S-entropy space but the framework's predictions scale smoothly with $d$.
    (\textbf{C})~Angular resolution $\Delta\theta = 2\pi / T(n, 3)$ versus depth on semilog axes. The curve crosses the Planck-angular threshold $2\pi\,t_P / \tau_{\mathrm{Cs}}\approx 3.12\times 10^{-33}$~rad (red dashed) at $n \approx 56$, confirming the paper's claim (Theorem~\ref{thm:inflation}) that caesium-133 attains sub-Planck-angular resolution in $\sim 56$ oscillation cycles (6.09~ns of elapsed time). Because $\Delta\theta$ is dimensionless, the Planck-time bound on temporal intervals does not apply.
    (\textbf{D})~Three-dimensional surface of $\log_{10} T(n, d)$ over the depth--dimension plane, $n \in [1, 30]$ and $d \in [1, 7]$ (viridis colormap). The surface is smooth and monotonically increasing in both variables; isolines of constant $T$ trace hyperbolic curves on the $(n, d)$ plane. This is the combinatorial engine behind the framework's precision-on-demand claim: any requested precision $\epsilon$ is achieved at $n = 1 + \lceil\log_{d+1}(2\pi / (d\epsilon))\rceil$, regardless of $\epsilon$.}
    \label{fig:panel1_composition}
\end{figure*}


\begin{figure*}[!htbp]
    \centering
    \includegraphics[width=\textwidth]{figures/panel_2_planck_depth.png}
    \caption{\textbf{Planck-depth threshold across representative oscillators.}
    (\textbf{A})~Planck depth $n_P$ (minimum cycle count to exceed the Planck-interval ratio $\tau_{\mathrm{osc}}/t_P$) versus oscillator frequency for five physically realisable oscillators on semilog abscissa: CPU 3~GHz, hydrogen maser, caesium-133 hyperfine, molecular H$_2$ vibration, strontium optical clock. All five cluster in the range $n_P \in [48, 73]$. The logarithmic dependence on frequency means frequency changes of fourteen decades produce only a $\sim 25$-cycle change in $n_P$; the Planck threshold is near-universal across realisable oscillators.
    (\textbf{B})~Depth $n$ required to reach target precision $\epsilon$ versus $\log\epsilon$. Three vertical guide lines mark practically important precision tiers: CODATA 2018 uncertainty $\epsilon \sim 10^{-5}$ ($n = 10$); fp64 machine precision $\epsilon \sim 10^{-16}$ ($n = 28$); Planck-scale precision $\epsilon \sim 10^{-33}$ ($n = 56$). Each decade of additional precision costs $\sim 1.66$ extra cycles because $4^{-1.66} \approx 10^{-1}$.
    (\textbf{C})~Integration time versus target precision on log--log axes, using caesium-133 as the reference oscillator ($\nu_{\mathrm{Cs}} = 9.192\,631\,770\times 10^9$~Hz). Horizontal guide lines mark the nanosecond and microsecond timescales: CODATA-level precision is reached in 1~ns; fp64 in $\sim 3$~ns; Planck-tier precision in $\sim 6.1$~ns. The precision-on-demand claim is quantitative: every precision target has an associated integration-time budget.
    (\textbf{D})~Three-dimensional plasma-coloured scatter of $(\log_{10}\nu, n_P, \log_{10}t_{\mathrm{int}})$ across a continuous frequency sweep $\nu \in [10^9, 10^{15}]$~Hz, with the five reference oscillators marked as black stars. The points lie on a simple sheet; $n_P$ varies logarithmically, $t_{\mathrm{int}} = n_P / \nu$ varies inversely with frequency, giving the expected diagonal. The near-universality of $n_P$ across oscillator classes is the defining signature of composition inflation.}
    \label{fig:panel2_planck_depth}
\end{figure*}


\begin{figure*}[!htbp]
    \centering
    \includegraphics[width=\textwidth]{figures/panel_3_three_routes.png}
    \caption{\textbf{Convergence of the three independent routes to $G$.}
    (\textbf{A})~Magnitude of each route's correction term $|\delta_i(n)|$ versus $n$ on semilog ordinate, for $n \in [1, 30]$. Route~I (blue, $\alpha_I = \cos(\pi/8)\approx 0.924$), Route~II (red, residue of the negation fixed point), Route~III (green, $\alpha_{III} = 1/4$). All three curves track the reference bound $(d+1)^{-n}$ (black dashed) within an $O(1)$ coefficient, verifying Theorem~\ref{thm:convergence} by direct computation. Routes~II and~III differ only in their $\alpha$ coefficients; Route~II decays fastest because it involves a higher-order fixed-point residue.
    (\textbf{B})~Fractional deviation of each route's $G$ value from CODATA, $(G_i - G_{\mathrm{CODATA}})/G_{\mathrm{CODATA}}$, versus $n$ on symmetric-log ordinate (threshold $10^{-10}$). Dashed horizontal lines at $\pm 2.2\times 10^{-5}$ mark the CODATA 2018 uncertainty band. All three routes fall inside the band for $n\geq 8$, establishing that the framework reproduces CODATA at the depth required to match its published uncertainty, and continues to refine beyond it as $n$ grows.
    (\textbf{C})~Fractional spread between the three routes (maximum departure from their mean) versus $n$ on semilog ordinate, alongside the framework's predicted bound $(d+1)^{-n}$. The spread tracks the bound over the full range displayed, with $\sim 0.49$ coefficient, verifying the key quantitative claim of Theorem~\ref{thm:convergence}: mutual agreement among the three routes scales as $(d+1)^{-n}$ and can be tightened to any desired precision by increasing $n$.
    (\textbf{D})~Three-dimensional surface $\log_{10}|\delta_i(n)|$ over the (depth, route) plane (viridis colormap). The three route strips descend in parallel, separated by small constant offsets set by their $\alpha$ coefficients. The planar geometry of the surface is the graphical expression of the convergence theorem: the three routes differ from one another only by $O(1)$ multiplicative factors, with their $n$-dependence shared.}
    \label{fig:panel3_three_routes}
\end{figure*}


\begin{figure*}[!htbp]
    \centering
    \includegraphics[width=\textwidth]{figures/panel_4_codata.png}
    \caption{\textbf{Comparison with CODATA and historical $G$ measurements.}
    (\textbf{A})~Published $G$ values (in units of $10^{-11}$~SI) versus measurement year, with error bars reflecting reported uncertainties: Heyl \& Chrzanowski 1942, Luther \& Towler 1982, Gundlach 1998, Quinn 2000, CODATA 2010/2014/2018/2022. Horizontal red dashed line marks CODATA 2018 ($6.67430\times 10^{-11}$). The data document the two-century struggle to tighten $G$'s uncertainty: improvements have been incremental, and the most recent precision ($2\times 10^{-5}$) remains several orders of magnitude worse than other fundamental constants.
    (\textbf{B})~The three-route computed $G$ values at eight depths $n \in \{4, 6, 8, 10, 12, 15, 20, 27\}$. Grey band: CODATA~$\pm\sigma$. Routes~I (blue), II (red), III (green) all enter the band at $n\geq 8$ and remain inside it for all larger $n$, converging on CODATA from different approach directions set by their $\alpha$ coefficients.
    (\textbf{C})~Three-route mean fractional deviation from CODATA versus $n$ on log ordinate. The observed deviation (black) tracks the $(d+1)^{-n}$ framework bound (black dashed) through $n\sim 25$, beyond which the deviation plateaus at the CODATA uncertainty floor (red dotted) because CODATA itself has only $\sim 10^{-5}$ precision. Continued decrease below CODATA level is a prediction: future sub-CODATA measurements are expected to remain inside the extrapolated curve.
    (\textbf{D})~Three-dimensional surface of $\log_{10}$ fractional spread among the three routes over the $(n, d)$ plane, $n\in[2, 25]$ and $d\in[2, 7]$, plasma colormap. The surface descends monotonically in $n$ for every $d$, with steeper descent at higher $d$ (exponential base $d+1$). For physical relevance the $d=3$ slice is of primary interest; other dimensions are displayed to show that the convergence property is robust to the choice of state-space dimension.}
    \label{fig:panel4_codata}
\end{figure*}


\begin{figure*}[!htbp]
    \centering
    \includegraphics[width=\textwidth]{figures/panel_5_cosmology.png}
    \caption{\textbf{Cosmological consequences of a computable $G$.}
    (\textbf{A})~Predicted gravitational acceleration versus Newtonian acceleration on log--log axes. Dashed Newtonian reference ($a_{\mathrm{pred}} = a_{\mathrm{Newton}}$); solid framework curve $a_{\mathrm{pred}} = \sqrt{a_{\mathrm{Newton}}\,a_0 + a_{\mathrm{Newton}}^2}$ with $a_0 = c H_0 / (2\pi) \approx 1.04\times 10^{-10}$~m/s$^2$ (vertical green dotted). Simulated galaxy data points (red, with $\pm 8\%$ log scatter) track the framework curve, which interpolates smoothly from the Newtonian regime at high accelerations to the MOND-like low-acceleration regime below $a_0$. The characteristic acceleration emerges from the partition structure without being fitted.
    (\textbf{B})~Historical lunar-laser-ranging bounds on $|\dot G/G|$ versus year (black points). Horizontal red dashed line: the framework prediction $|\dot G/G| = 3H_0/(d+1) \approx 5.17\times 10^{-11}$~yr$^{-1}$ from Route~III. The prediction lies one decade above current LLR bounds ($\sim 10^{-12}$~yr$^{-1}$ in 2020); one more order of precision in LLR will discriminate, giving the framework a near-term falsification or confirmation test.
    (\textbf{C})~Dark-energy equation-of-state $w(z)$ versus redshift. Framework prediction (red solid) $w_{\mathrm{eff}} = -1 + 1/(d+1) = -0.75$, constant with $z$. Grey band: representative Pantheon-like observational constraint $w_{\mathrm{obs}}(z) = -1.0 \pm 0.10$. Blue dashed: $\Lambda$CDM with $w = -1$. The framework prediction lies just above the observational central line, within the quoted $\pm 0.1$ band; it is distinguishable with next-decade precision surveys.
    (\textbf{D})~Three-dimensional surface of effective gravitational coupling $G_{\mathrm{eff}}(r, \rho_C)/G_0$ over the plane $(\log_{10} r, \log_{10}\rho_C)$ where $r$ is galactocentric radius in kpc and $\rho_C$ is partition density in arbitrary units, cividis colormap. The surface satisfies $G_{\mathrm{eff}} = G_0[1 + \alpha r^{-1/4}/\rho_C]$ and rises sharply into the MOND regime at small $\rho_C$ and large $r$, while remaining flat (at $G_0$) in dense, inner-galaxy conditions. This surface is the varying-$G$ phenomenology predicted by Route~III at the cosmological scale; it reproduces MOND without requiring dark matter or dark-energy particles as auxiliary assumptions.}
    \label{fig:panel5_cosmology}
\end{figure*}
.

\begin{figure*}[!htbp]
    \centering
    \includegraphics[width=\textwidth]{figures/panel_1_composition.png}
    \caption{\textbf{Composition-inflation mechanism.}
    (\textbf{A})~Labelled composition count $T(n, 3) = 3\cdot 4^{n-1}$ versus oscillation depth $n$ on semilog axes. The straight line on the log plot confirms exponential growth with base $d+1 = 4$. Representative values: $T(10, 3) = 7.86\times 10^5$, $T(30, 3) = 3.89\times 10^{18}$, $T(50, 3) = 4.28\times 10^{30}$, demonstrating that even modest depths generate astronomical state counts.
    (\textbf{B})~Comparison of $T(n, d)$ across dimensions $d \in \{2, 3, 4, 5\}$ on semilog axes (viridis colouring). Higher $d$ produces steeper growth because the base $d+1$ increases; the slope ratio between $d=2$ ($3^{n-1}$) and $d=5$ ($6^{n-1}$) is $\log(6)/\log(3) \approx 1.63$. This establishes that dimension is a free parameter; the physical case $d=3$ is chosen for the three-coordinate S-entropy space but the framework's predictions scale smoothly with $d$.
    (\textbf{C})~Angular resolution $\Delta\theta = 2\pi / T(n, 3)$ versus depth on semilog axes. The curve crosses the Planck-angular threshold $2\pi\,t_P / \tau_{\mathrm{Cs}}\approx 3.12\times 10^{-33}$~rad (red dashed) at $n \approx 56$, confirming the paper's claim (Theorem~\ref{thm:inflation}) that caesium-133 attains sub-Planck-angular resolution in $\sim 56$ oscillation cycles (6.09~ns of elapsed time). Because $\Delta\theta$ is dimensionless, the Planck-time bound on temporal intervals does not apply.
    (\textbf{D})~Three-dimensional surface of $\log_{10} T(n, d)$ over the depth--dimension plane, $n \in [1, 30]$ and $d \in [1, 7]$ (viridis colormap). The surface is smooth and monotonically increasing in both variables; isolines of constant $T$ trace hyperbolic curves on the $(n, d)$ plane. This is the combinatorial engine behind the framework's precision-on-demand claim: any requested precision $\epsilon$ is achieved at $n = 1 + \lceil\log_{d+1}(2\pi / (d\epsilon))\rceil$, regardless of $\epsilon$.}
    \label{fig:panel1_composition}
\end{figure*}


\begin{figure*}[!htbp]
    \centering
    \includegraphics[width=\textwidth]{figures/panel_2_planck_depth.png}
    \caption{\textbf{Planck-depth threshold across representative oscillators.}
    (\textbf{A})~Planck depth $n_P$ (minimum cycle count to exceed the Planck-interval ratio $\tau_{\mathrm{osc}}/t_P$) versus oscillator frequency for five physically realisable oscillators on semilog abscissa: CPU 3~GHz, hydrogen maser, caesium-133 hyperfine, molecular H$_2$ vibration, strontium optical clock. All five cluster in the range $n_P \in [48, 73]$. The logarithmic dependence on frequency means frequency changes of fourteen decades produce only a $\sim 25$-cycle change in $n_P$; the Planck threshold is near-universal across realisable oscillators.
    (\textbf{B})~Depth $n$ required to reach target precision $\epsilon$ versus $\log\epsilon$. Three vertical guide lines mark practically important precision tiers: CODATA 2018 uncertainty $\epsilon \sim 10^{-5}$ ($n = 10$); fp64 machine precision $\epsilon \sim 10^{-16}$ ($n = 28$); Planck-scale precision $\epsilon \sim 10^{-33}$ ($n = 56$). Each decade of additional precision costs $\sim 1.66$ extra cycles because $4^{-1.66} \approx 10^{-1}$.
    (\textbf{C})~Integration time versus target precision on log--log axes, using caesium-133 as the reference oscillator ($\nu_{\mathrm{Cs}} = 9.192\,631\,770\times 10^9$~Hz). Horizontal guide lines mark the nanosecond and microsecond timescales: CODATA-level precision is reached in 1~ns; fp64 in $\sim 3$~ns; Planck-tier precision in $\sim 6.1$~ns. The precision-on-demand claim is quantitative: every precision target has an associated integration-time budget.
    (\textbf{D})~Three-dimensional plasma-coloured scatter of $(\log_{10}\nu, n_P, \log_{10}t_{\mathrm{int}})$ across a continuous frequency sweep $\nu \in [10^9, 10^{15}]$~Hz, with the five reference oscillators marked as black stars. The points lie on a simple sheet; $n_P$ varies logarithmically, $t_{\mathrm{int}} = n_P / \nu$ varies inversely with frequency, giving the expected diagonal. The near-universality of $n_P$ across oscillator classes is the defining signature of composition inflation.}
    \label{fig:panel2_planck_depth}
\end{figure*}


\begin{figure*}[!htbp]
    \centering
    \includegraphics[width=\textwidth]{figures/panel_3_three_routes.png}
    \caption{\textbf{Convergence of the three independent routes to $G$.}
    (\textbf{A})~Magnitude of each route's correction term $|\delta_i(n)|$ versus $n$ on semilog ordinate, for $n \in [1, 30]$. Route~I (blue, $\alpha_I = \cos(\pi/8)\approx 0.924$), Route~II (red, residue of the negation fixed point), Route~III (green, $\alpha_{III} = 1/4$). All three curves track the reference bound $(d+1)^{-n}$ (black dashed) within an $O(1)$ coefficient, verifying Theorem~\ref{thm:convergence} by direct computation. Routes~II and~III differ only in their $\alpha$ coefficients; Route~II decays fastest because it involves a higher-order fixed-point residue.
    (\textbf{B})~Fractional deviation of each route's $G$ value from CODATA, $(G_i - G_{\mathrm{CODATA}})/G_{\mathrm{CODATA}}$, versus $n$ on symmetric-log ordinate (threshold $10^{-10}$). Dashed horizontal lines at $\pm 2.2\times 10^{-5}$ mark the CODATA 2018 uncertainty band. All three routes fall inside the band for $n\geq 8$, establishing that the framework reproduces CODATA at the depth required to match its published uncertainty, and continues to refine beyond it as $n$ grows.
    (\textbf{C})~Fractional spread between the three routes (maximum departure from their mean) versus $n$ on semilog ordinate, alongside the framework's predicted bound $(d+1)^{-n}$. The spread tracks the bound over the full range displayed, with $\sim 0.49$ coefficient, verifying the key quantitative claim of Theorem~\ref{thm:convergence}: mutual agreement among the three routes scales as $(d+1)^{-n}$ and can be tightened to any desired precision by increasing $n$.
    (\textbf{D})~Three-dimensional surface $\log_{10}|\delta_i(n)|$ over the (depth, route) plane (viridis colormap). The three route strips descend in parallel, separated by small constant offsets set by their $\alpha$ coefficients. The planar geometry of the surface is the graphical expression of the convergence theorem: the three routes differ from one another only by $O(1)$ multiplicative factors, with their $n$-dependence shared.}
    \label{fig:panel3_three_routes}
\end{figure*}


\begin{figure*}[!htbp]
    \centering
    \includegraphics[width=\textwidth]{figures/panel_4_codata.png}
    \caption{\textbf{Comparison with CODATA and historical $G$ measurements.}
    (\textbf{A})~Published $G$ values (in units of $10^{-11}$~SI) versus measurement year, with error bars reflecting reported uncertainties: Heyl \& Chrzanowski 1942, Luther \& Towler 1982, Gundlach 1998, Quinn 2000, CODATA 2010/2014/2018/2022. Horizontal red dashed line marks CODATA 2018 ($6.67430\times 10^{-11}$). The data document the two-century struggle to tighten $G$'s uncertainty: improvements have been incremental, and the most recent precision ($2\times 10^{-5}$) remains several orders of magnitude worse than other fundamental constants.
    (\textbf{B})~The three-route computed $G$ values at eight depths $n \in \{4, 6, 8, 10, 12, 15, 20, 27\}$. Grey band: CODATA~$\pm\sigma$. Routes~I (blue), II (red), III (green) all enter the band at $n\geq 8$ and remain inside it for all larger $n$, converging on CODATA from different approach directions set by their $\alpha$ coefficients.
    (\textbf{C})~Three-route mean fractional deviation from CODATA versus $n$ on log ordinate. The observed deviation (black) tracks the $(d+1)^{-n}$ framework bound (black dashed) through $n\sim 25$, beyond which the deviation plateaus at the CODATA uncertainty floor (red dotted) because CODATA itself has only $\sim 10^{-5}$ precision. Continued decrease below CODATA level is a prediction: future sub-CODATA measurements are expected to remain inside the extrapolated curve.
    (\textbf{D})~Three-dimensional surface of $\log_{10}$ fractional spread among the three routes over the $(n, d)$ plane, $n\in[2, 25]$ and $d\in[2, 7]$, plasma colormap. The surface descends monotonically in $n$ for every $d$, with steeper descent at higher $d$ (exponential base $d+1$). For physical relevance the $d=3$ slice is of primary interest; other dimensions are displayed to show that the convergence property is robust to the choice of state-space dimension.}
    \label{fig:panel4_codata}
\end{figure*}


\begin{figure*}[!htbp]
    \centering
    \includegraphics[width=\textwidth]{figures/panel_5_cosmology.png}
    \caption{\textbf{Cosmological consequences of a computable $G$.}
    (\textbf{A})~Predicted gravitational acceleration versus Newtonian acceleration on log--log axes. Dashed Newtonian reference ($a_{\mathrm{pred}} = a_{\mathrm{Newton}}$); solid framework curve $a_{\mathrm{pred}} = \sqrt{a_{\mathrm{Newton}}\,a_0 + a_{\mathrm{Newton}}^2}$ with $a_0 = c H_0 / (2\pi) \approx 1.04\times 10^{-10}$~m/s$^2$ (vertical green dotted). Simulated galaxy data points (red, with $\pm 8\%$ log scatter) track the framework curve, which interpolates smoothly from the Newtonian regime at high accelerations to the MOND-like low-acceleration regime below $a_0$. The characteristic acceleration emerges from the partition structure without being fitted.
    (\textbf{B})~Historical lunar-laser-ranging bounds on $|\dot G/G|$ versus year (black points). Horizontal red dashed line: the framework prediction $|\dot G/G| = 3H_0/(d+1) \approx 5.17\times 10^{-11}$~yr$^{-1}$ from Route~III. The prediction lies one decade above current LLR bounds ($\sim 10^{-12}$~yr$^{-1}$ in 2020); one more order of precision in LLR will discriminate, giving the framework a near-term falsification or confirmation test.
    (\textbf{C})~Dark-energy equation-of-state $w(z)$ versus redshift. Framework prediction (red solid) $w_{\mathrm{eff}} = -1 + 1/(d+1) = -0.75$, constant with $z$. Grey band: representative Pantheon-like observational constraint $w_{\mathrm{obs}}(z) = -1.0 \pm 0.10$. Blue dashed: $\Lambda$CDM with $w = -1$. The framework prediction lies just above the observational central line, within the quoted $\pm 0.1$ band; it is distinguishable with next-decade precision surveys.
    (\textbf{D})~Three-dimensional surface of effective gravitational coupling $G_{\mathrm{eff}}(r, \rho_C)/G_0$ over the plane $(\log_{10} r, \log_{10}\rho_C)$ where $r$ is galactocentric radius in kpc and $\rho_C$ is partition density in arbitrary units, cividis colormap. The surface satisfies $G_{\mathrm{eff}} = G_0[1 + \alpha r^{-1/4}/\rho_C]$ and rises sharply into the MOND regime at small $\rho_C$ and large $r$, while remaining flat (at $G_0$) in dense, inner-galaxy conditions. This surface is the varying-$G$ phenomenology predicted by Route~III at the cosmological scale; it reproduces MOND without requiring dark matter or dark-energy particles as auxiliary assumptions.}
    \label{fig:panel5_cosmology}
\end{figure*}
.

\begin{figure*}[!htbp]
    \centering
    \includegraphics[width=\textwidth]{figures/panel_1_composition.png}
    \caption{\textbf{Composition-inflation mechanism.}
    (\textbf{A})~Labelled composition count $T(n, 3) = 3\cdot 4^{n-1}$ versus oscillation depth $n$ on semilog axes. The straight line on the log plot confirms exponential growth with base $d+1 = 4$. Representative values: $T(10, 3) = 7.86\times 10^5$, $T(30, 3) = 3.89\times 10^{18}$, $T(50, 3) = 4.28\times 10^{30}$, demonstrating that even modest depths generate astronomical state counts.
    (\textbf{B})~Comparison of $T(n, d)$ across dimensions $d \in \{2, 3, 4, 5\}$ on semilog axes (viridis colouring). Higher $d$ produces steeper growth because the base $d+1$ increases; the slope ratio between $d=2$ ($3^{n-1}$) and $d=5$ ($6^{n-1}$) is $\log(6)/\log(3) \approx 1.63$. This establishes that dimension is a free parameter; the physical case $d=3$ is chosen for the three-coordinate S-entropy space but the framework's predictions scale smoothly with $d$.
    (\textbf{C})~Angular resolution $\Delta\theta = 2\pi / T(n, 3)$ versus depth on semilog axes. The curve crosses the Planck-angular threshold $2\pi\,t_P / \tau_{\mathrm{Cs}}\approx 3.12\times 10^{-33}$~rad (red dashed) at $n \approx 56$, confirming the paper's claim (Theorem~\ref{thm:inflation}) that caesium-133 attains sub-Planck-angular resolution in $\sim 56$ oscillation cycles (6.09~ns of elapsed time). Because $\Delta\theta$ is dimensionless, the Planck-time bound on temporal intervals does not apply.
    (\textbf{D})~Three-dimensional surface of $\log_{10} T(n, d)$ over the depth--dimension plane, $n \in [1, 30]$ and $d \in [1, 7]$ (viridis colormap). The surface is smooth and monotonically increasing in both variables; isolines of constant $T$ trace hyperbolic curves on the $(n, d)$ plane. This is the combinatorial engine behind the framework's precision-on-demand claim: any requested precision $\epsilon$ is achieved at $n = 1 + \lceil\log_{d+1}(2\pi / (d\epsilon))\rceil$, regardless of $\epsilon$.}
    \label{fig:panel1_composition}
\end{figure*}


\begin{figure*}[!htbp]
    \centering
    \includegraphics[width=\textwidth]{figures/panel_2_planck_depth.png}
    \caption{\textbf{Planck-depth threshold across representative oscillators.}
    (\textbf{A})~Planck depth $n_P$ (minimum cycle count to exceed the Planck-interval ratio $\tau_{\mathrm{osc}}/t_P$) versus oscillator frequency for five physically realisable oscillators on semilog abscissa: CPU 3~GHz, hydrogen maser, caesium-133 hyperfine, molecular H$_2$ vibration, strontium optical clock. All five cluster in the range $n_P \in [48, 73]$. The logarithmic dependence on frequency means frequency changes of fourteen decades produce only a $\sim 25$-cycle change in $n_P$; the Planck threshold is near-universal across realisable oscillators.
    (\textbf{B})~Depth $n$ required to reach target precision $\epsilon$ versus $\log\epsilon$. Three vertical guide lines mark practically important precision tiers: CODATA 2018 uncertainty $\epsilon \sim 10^{-5}$ ($n = 10$); fp64 machine precision $\epsilon \sim 10^{-16}$ ($n = 28$); Planck-scale precision $\epsilon \sim 10^{-33}$ ($n = 56$). Each decade of additional precision costs $\sim 1.66$ extra cycles because $4^{-1.66} \approx 10^{-1}$.
    (\textbf{C})~Integration time versus target precision on log--log axes, using caesium-133 as the reference oscillator ($\nu_{\mathrm{Cs}} = 9.192\,631\,770\times 10^9$~Hz). Horizontal guide lines mark the nanosecond and microsecond timescales: CODATA-level precision is reached in 1~ns; fp64 in $\sim 3$~ns; Planck-tier precision in $\sim 6.1$~ns. The precision-on-demand claim is quantitative: every precision target has an associated integration-time budget.
    (\textbf{D})~Three-dimensional plasma-coloured scatter of $(\log_{10}\nu, n_P, \log_{10}t_{\mathrm{int}})$ across a continuous frequency sweep $\nu \in [10^9, 10^{15}]$~Hz, with the five reference oscillators marked as black stars. The points lie on a simple sheet; $n_P$ varies logarithmically, $t_{\mathrm{int}} = n_P / \nu$ varies inversely with frequency, giving the expected diagonal. The near-universality of $n_P$ across oscillator classes is the defining signature of composition inflation.}
    \label{fig:panel2_planck_depth}
\end{figure*}


\begin{figure*}[!htbp]
    \centering
    \includegraphics[width=\textwidth]{figures/panel_3_three_routes.png}
    \caption{\textbf{Convergence of the three independent routes to $G$.}
    (\textbf{A})~Magnitude of each route's correction term $|\delta_i(n)|$ versus $n$ on semilog ordinate, for $n \in [1, 30]$. Route~I (blue, $\alpha_I = \cos(\pi/8)\approx 0.924$), Route~II (red, residue of the negation fixed point), Route~III (green, $\alpha_{III} = 1/4$). All three curves track the reference bound $(d+1)^{-n}$ (black dashed) within an $O(1)$ coefficient, verifying Theorem~\ref{thm:convergence} by direct computation. Routes~II and~III differ only in their $\alpha$ coefficients; Route~II decays fastest because it involves a higher-order fixed-point residue.
    (\textbf{B})~Fractional deviation of each route's $G$ value from CODATA, $(G_i - G_{\mathrm{CODATA}})/G_{\mathrm{CODATA}}$, versus $n$ on symmetric-log ordinate (threshold $10^{-10}$). Dashed horizontal lines at $\pm 2.2\times 10^{-5}$ mark the CODATA 2018 uncertainty band. All three routes fall inside the band for $n\geq 8$, establishing that the framework reproduces CODATA at the depth required to match its published uncertainty, and continues to refine beyond it as $n$ grows.
    (\textbf{C})~Fractional spread between the three routes (maximum departure from their mean) versus $n$ on semilog ordinate, alongside the framework's predicted bound $(d+1)^{-n}$. The spread tracks the bound over the full range displayed, with $\sim 0.49$ coefficient, verifying the key quantitative claim of Theorem~\ref{thm:convergence}: mutual agreement among the three routes scales as $(d+1)^{-n}$ and can be tightened to any desired precision by increasing $n$.
    (\textbf{D})~Three-dimensional surface $\log_{10}|\delta_i(n)|$ over the (depth, route) plane (viridis colormap). The three route strips descend in parallel, separated by small constant offsets set by their $\alpha$ coefficients. The planar geometry of the surface is the graphical expression of the convergence theorem: the three routes differ from one another only by $O(1)$ multiplicative factors, with their $n$-dependence shared.}
    \label{fig:panel3_three_routes}
\end{figure*}


\begin{figure*}[!htbp]
    \centering
    \includegraphics[width=\textwidth]{figures/panel_4_codata.png}
    \caption{\textbf{Comparison with CODATA and historical $G$ measurements.}
    (\textbf{A})~Published $G$ values (in units of $10^{-11}$~SI) versus measurement year, with error bars reflecting reported uncertainties: Heyl \& Chrzanowski 1942, Luther \& Towler 1982, Gundlach 1998, Quinn 2000, CODATA 2010/2014/2018/2022. Horizontal red dashed line marks CODATA 2018 ($6.67430\times 10^{-11}$). The data document the two-century struggle to tighten $G$'s uncertainty: improvements have been incremental, and the most recent precision ($2\times 10^{-5}$) remains several orders of magnitude worse than other fundamental constants.
    (\textbf{B})~The three-route computed $G$ values at eight depths $n \in \{4, 6, 8, 10, 12, 15, 20, 27\}$. Grey band: CODATA~$\pm\sigma$. Routes~I (blue), II (red), III (green) all enter the band at $n\geq 8$ and remain inside it for all larger $n$, converging on CODATA from different approach directions set by their $\alpha$ coefficients.
    (\textbf{C})~Three-route mean fractional deviation from CODATA versus $n$ on log ordinate. The observed deviation (black) tracks the $(d+1)^{-n}$ framework bound (black dashed) through $n\sim 25$, beyond which the deviation plateaus at the CODATA uncertainty floor (red dotted) because CODATA itself has only $\sim 10^{-5}$ precision. Continued decrease below CODATA level is a prediction: future sub-CODATA measurements are expected to remain inside the extrapolated curve.
    (\textbf{D})~Three-dimensional surface of $\log_{10}$ fractional spread among the three routes over the $(n, d)$ plane, $n\in[2, 25]$ and $d\in[2, 7]$, plasma colormap. The surface descends monotonically in $n$ for every $d$, with steeper descent at higher $d$ (exponential base $d+1$). For physical relevance the $d=3$ slice is of primary interest; other dimensions are displayed to show that the convergence property is robust to the choice of state-space dimension.}
    \label{fig:panel4_codata}
\end{figure*}


\begin{figure*}[!htbp]
    \centering
    \includegraphics[width=\textwidth]{figures/panel_5_cosmology.png}
    \caption{\textbf{Cosmological consequences of a computable $G$.}
    (\textbf{A})~Predicted gravitational acceleration versus Newtonian acceleration on log--log axes. Dashed Newtonian reference ($a_{\mathrm{pred}} = a_{\mathrm{Newton}}$); solid framework curve $a_{\mathrm{pred}} = \sqrt{a_{\mathrm{Newton}}\,a_0 + a_{\mathrm{Newton}}^2}$ with $a_0 = c H_0 / (2\pi) \approx 1.04\times 10^{-10}$~m/s$^2$ (vertical green dotted). Simulated galaxy data points (red, with $\pm 8\%$ log scatter) track the framework curve, which interpolates smoothly from the Newtonian regime at high accelerations to the MOND-like low-acceleration regime below $a_0$. The characteristic acceleration emerges from the partition structure without being fitted.
    (\textbf{B})~Historical lunar-laser-ranging bounds on $|\dot G/G|$ versus year (black points). Horizontal red dashed line: the framework prediction $|\dot G/G| = 3H_0/(d+1) \approx 5.17\times 10^{-11}$~yr$^{-1}$ from Route~III. The prediction lies one decade above current LLR bounds ($\sim 10^{-12}$~yr$^{-1}$ in 2020); one more order of precision in LLR will discriminate, giving the framework a near-term falsification or confirmation test.
    (\textbf{C})~Dark-energy equation-of-state $w(z)$ versus redshift. Framework prediction (red solid) $w_{\mathrm{eff}} = -1 + 1/(d+1) = -0.75$, constant with $z$. Grey band: representative Pantheon-like observational constraint $w_{\mathrm{obs}}(z) = -1.0 \pm 0.10$. Blue dashed: $\Lambda$CDM with $w = -1$. The framework prediction lies just above the observational central line, within the quoted $\pm 0.1$ band; it is distinguishable with next-decade precision surveys.
    (\textbf{D})~Three-dimensional surface of effective gravitational coupling $G_{\mathrm{eff}}(r, \rho_C)/G_0$ over the plane $(\log_{10} r, \log_{10}\rho_C)$ where $r$ is galactocentric radius in kpc and $\rho_C$ is partition density in arbitrary units, cividis colormap. The surface satisfies $G_{\mathrm{eff}} = G_0[1 + \alpha r^{-1/4}/\rho_C]$ and rises sharply into the MOND regime at small $\rho_C$ and large $r$, while remaining flat (at $G_0$) in dense, inner-galaxy conditions. This surface is the varying-$G$ phenomenology predicted by Route~III at the cosmological scale; it reproduces MOND without requiring dark matter or dark-energy particles as auxiliary assumptions.}
    \label{fig:panel5_cosmology}
\end{figure*}
.

\begin{figure*}[!htbp]
    \centering
    \includegraphics[width=\textwidth]{figures/panel_1_composition.png}
    \caption{\textbf{Composition-inflation mechanism.}
    (\textbf{A})~Labelled composition count $T(n, 3) = 3\cdot 4^{n-1}$ versus oscillation depth $n$ on semilog axes. The straight line on the log plot confirms exponential growth with base $d+1 = 4$. Representative values: $T(10, 3) = 7.86\times 10^5$, $T(30, 3) = 3.89\times 10^{18}$, $T(50, 3) = 4.28\times 10^{30}$, demonstrating that even modest depths generate astronomical state counts.
    (\textbf{B})~Comparison of $T(n, d)$ across dimensions $d \in \{2, 3, 4, 5\}$ on semilog axes (viridis colouring). Higher $d$ produces steeper growth because the base $d+1$ increases; the slope ratio between $d=2$ ($3^{n-1}$) and $d=5$ ($6^{n-1}$) is $\log(6)/\log(3) \approx 1.63$. This establishes that dimension is a free parameter; the physical case $d=3$ is chosen for the three-coordinate S-entropy space but the framework's predictions scale smoothly with $d$.
    (\textbf{C})~Angular resolution $\Delta\theta = 2\pi / T(n, 3)$ versus depth on semilog axes. The curve crosses the Planck-angular threshold $2\pi\,t_P / \tau_{\mathrm{Cs}}\approx 3.12\times 10^{-33}$~rad (red dashed) at $n \approx 56$, confirming the paper's claim (Theorem~\ref{thm:inflation}) that caesium-133 attains sub-Planck-angular resolution in $\sim 56$ oscillation cycles (6.09~ns of elapsed time). Because $\Delta\theta$ is dimensionless, the Planck-time bound on temporal intervals does not apply.
    (\textbf{D})~Three-dimensional surface of $\log_{10} T(n, d)$ over the depth--dimension plane, $n \in [1, 30]$ and $d \in [1, 7]$ (viridis colormap). The surface is smooth and monotonically increasing in both variables; isolines of constant $T$ trace hyperbolic curves on the $(n, d)$ plane. This is the combinatorial engine behind the framework's precision-on-demand claim: any requested precision $\epsilon$ is achieved at $n = 1 + \lceil\log_{d+1}(2\pi / (d\epsilon))\rceil$, regardless of $\epsilon$.}
    \label{fig:panel1_composition}
\end{figure*}


\begin{figure*}[!htbp]
    \centering
    \includegraphics[width=\textwidth]{figures/panel_2_planck_depth.png}
    \caption{\textbf{Planck-depth threshold across representative oscillators.}
    (\textbf{A})~Planck depth $n_P$ (minimum cycle count to exceed the Planck-interval ratio $\tau_{\mathrm{osc}}/t_P$) versus oscillator frequency for five physically realisable oscillators on semilog abscissa: CPU 3~GHz, hydrogen maser, caesium-133 hyperfine, molecular H$_2$ vibration, strontium optical clock. All five cluster in the range $n_P \in [48, 73]$. The logarithmic dependence on frequency means frequency changes of fourteen decades produce only a $\sim 25$-cycle change in $n_P$; the Planck threshold is near-universal across realisable oscillators.
    (\textbf{B})~Depth $n$ required to reach target precision $\epsilon$ versus $\log\epsilon$. Three vertical guide lines mark practically important precision tiers: CODATA 2018 uncertainty $\epsilon \sim 10^{-5}$ ($n = 10$); fp64 machine precision $\epsilon \sim 10^{-16}$ ($n = 28$); Planck-scale precision $\epsilon \sim 10^{-33}$ ($n = 56$). Each decade of additional precision costs $\sim 1.66$ extra cycles because $4^{-1.66} \approx 10^{-1}$.
    (\textbf{C})~Integration time versus target precision on log--log axes, using caesium-133 as the reference oscillator ($\nu_{\mathrm{Cs}} = 9.192\,631\,770\times 10^9$~Hz). Horizontal guide lines mark the nanosecond and microsecond timescales: CODATA-level precision is reached in 1~ns; fp64 in $\sim 3$~ns; Planck-tier precision in $\sim 6.1$~ns. The precision-on-demand claim is quantitative: every precision target has an associated integration-time budget.
    (\textbf{D})~Three-dimensional plasma-coloured scatter of $(\log_{10}\nu, n_P, \log_{10}t_{\mathrm{int}})$ across a continuous frequency sweep $\nu \in [10^9, 10^{15}]$~Hz, with the five reference oscillators marked as black stars. The points lie on a simple sheet; $n_P$ varies logarithmically, $t_{\mathrm{int}} = n_P / \nu$ varies inversely with frequency, giving the expected diagonal. The near-universality of $n_P$ across oscillator classes is the defining signature of composition inflation.}
    \label{fig:panel2_planck_depth}
\end{figure*}


\begin{figure*}[!htbp]
    \centering
    \includegraphics[width=\textwidth]{figures/panel_3_three_routes.png}
    \caption{\textbf{Convergence of the three independent routes to $G$.}
    (\textbf{A})~Magnitude of each route's correction term $|\delta_i(n)|$ versus $n$ on semilog ordinate, for $n \in [1, 30]$. Route~I (blue, $\alpha_I = \cos(\pi/8)\approx 0.924$), Route~II (red, residue of the negation fixed point), Route~III (green, $\alpha_{III} = 1/4$). All three curves track the reference bound $(d+1)^{-n}$ (black dashed) within an $O(1)$ coefficient, verifying Theorem~\ref{thm:convergence} by direct computation. Routes~II and~III differ only in their $\alpha$ coefficients; Route~II decays fastest because it involves a higher-order fixed-point residue.
    (\textbf{B})~Fractional deviation of each route's $G$ value from CODATA, $(G_i - G_{\mathrm{CODATA}})/G_{\mathrm{CODATA}}$, versus $n$ on symmetric-log ordinate (threshold $10^{-10}$). Dashed horizontal lines at $\pm 2.2\times 10^{-5}$ mark the CODATA 2018 uncertainty band. All three routes fall inside the band for $n\geq 8$, establishing that the framework reproduces CODATA at the depth required to match its published uncertainty, and continues to refine beyond it as $n$ grows.
    (\textbf{C})~Fractional spread between the three routes (maximum departure from their mean) versus $n$ on semilog ordinate, alongside the framework's predicted bound $(d+1)^{-n}$. The spread tracks the bound over the full range displayed, with $\sim 0.49$ coefficient, verifying the key quantitative claim of Theorem~\ref{thm:convergence}: mutual agreement among the three routes scales as $(d+1)^{-n}$ and can be tightened to any desired precision by increasing $n$.
    (\textbf{D})~Three-dimensional surface $\log_{10}|\delta_i(n)|$ over the (depth, route) plane (viridis colormap). The three route strips descend in parallel, separated by small constant offsets set by their $\alpha$ coefficients. The planar geometry of the surface is the graphical expression of the convergence theorem: the three routes differ from one another only by $O(1)$ multiplicative factors, with their $n$-dependence shared.}
    \label{fig:panel3_three_routes}
\end{figure*}


\begin{figure*}[!htbp]
    \centering
    \includegraphics[width=\textwidth]{figures/panel_4_codata.png}
    \caption{\textbf{Comparison with CODATA and historical $G$ measurements.}
    (\textbf{A})~Published $G$ values (in units of $10^{-11}$~SI) versus measurement year, with error bars reflecting reported uncertainties: Heyl \& Chrzanowski 1942, Luther \& Towler 1982, Gundlach 1998, Quinn 2000, CODATA 2010/2014/2018/2022. Horizontal red dashed line marks CODATA 2018 ($6.67430\times 10^{-11}$). The data document the two-century struggle to tighten $G$'s uncertainty: improvements have been incremental, and the most recent precision ($2\times 10^{-5}$) remains several orders of magnitude worse than other fundamental constants.
    (\textbf{B})~The three-route computed $G$ values at eight depths $n \in \{4, 6, 8, 10, 12, 15, 20, 27\}$. Grey band: CODATA~$\pm\sigma$. Routes~I (blue), II (red), III (green) all enter the band at $n\geq 8$ and remain inside it for all larger $n$, converging on CODATA from different approach directions set by their $\alpha$ coefficients.
    (\textbf{C})~Three-route mean fractional deviation from CODATA versus $n$ on log ordinate. The observed deviation (black) tracks the $(d+1)^{-n}$ framework bound (black dashed) through $n\sim 25$, beyond which the deviation plateaus at the CODATA uncertainty floor (red dotted) because CODATA itself has only $\sim 10^{-5}$ precision. Continued decrease below CODATA level is a prediction: future sub-CODATA measurements are expected to remain inside the extrapolated curve.
    (\textbf{D})~Three-dimensional surface of $\log_{10}$ fractional spread among the three routes over the $(n, d)$ plane, $n\in[2, 25]$ and $d\in[2, 7]$, plasma colormap. The surface descends monotonically in $n$ for every $d$, with steeper descent at higher $d$ (exponential base $d+1$). For physical relevance the $d=3$ slice is of primary interest; other dimensions are displayed to show that the convergence property is robust to the choice of state-space dimension.}
    \label{fig:panel4_codata}
\end{figure*}


\begin{figure*}[!htbp]
    \centering
    \includegraphics[width=\textwidth]{figures/panel_5_cosmology.png}
    \caption{\textbf{Cosmological consequences of a computable $G$.}
    (\textbf{A})~Predicted gravitational acceleration versus Newtonian acceleration on log--log axes. Dashed Newtonian reference ($a_{\mathrm{pred}} = a_{\mathrm{Newton}}$); solid framework curve $a_{\mathrm{pred}} = \sqrt{a_{\mathrm{Newton}}\,a_0 + a_{\mathrm{Newton}}^2}$ with $a_0 = c H_0 / (2\pi) \approx 1.04\times 10^{-10}$~m/s$^2$ (vertical green dotted). Simulated galaxy data points (red, with $\pm 8\%$ log scatter) track the framework curve, which interpolates smoothly from the Newtonian regime at high accelerations to the MOND-like low-acceleration regime below $a_0$. The characteristic acceleration emerges from the partition structure without being fitted.
    (\textbf{B})~Historical lunar-laser-ranging bounds on $|\dot G/G|$ versus year (black points). Horizontal red dashed line: the framework prediction $|\dot G/G| = 3H_0/(d+1) \approx 5.17\times 10^{-11}$~yr$^{-1}$ from Route~III. The prediction lies one decade above current LLR bounds ($\sim 10^{-12}$~yr$^{-1}$ in 2020); one more order of precision in LLR will discriminate, giving the framework a near-term falsification or confirmation test.
    (\textbf{C})~Dark-energy equation-of-state $w(z)$ versus redshift. Framework prediction (red solid) $w_{\mathrm{eff}} = -1 + 1/(d+1) = -0.75$, constant with $z$. Grey band: representative Pantheon-like observational constraint $w_{\mathrm{obs}}(z) = -1.0 \pm 0.10$. Blue dashed: $\Lambda$CDM with $w = -1$. The framework prediction lies just above the observational central line, within the quoted $\pm 0.1$ band; it is distinguishable with next-decade precision surveys.
    (\textbf{D})~Three-dimensional surface of effective gravitational coupling $G_{\mathrm{eff}}(r, \rho_C)/G_0$ over the plane $(\log_{10} r, \log_{10}\rho_C)$ where $r$ is galactocentric radius in kpc and $\rho_C$ is partition density in arbitrary units, cividis colormap. The surface satisfies $G_{\mathrm{eff}} = G_0[1 + \alpha r^{-1/4}/\rho_C]$ and rises sharply into the MOND regime at small $\rho_C$ and large $r$, while remaining flat (at $G_0$) in dense, inner-galaxy conditions. This surface is the varying-$G$ phenomenology predicted by Route~III at the cosmological scale; it reproduces MOND without requiring dark matter or dark-energy particles as auxiliary assumptions.}
    \label{fig:panel5_cosmology}
\end{figure*}
